% Part: first-order-logic
% Chapter: natural-deduction
% Section: rules-and-proofs

\documentclass[../../../include/open-logic-section]{subfiles}

\begin{document}

\iftag{FOL}
      {\olfileid{fol}{ntd}{rul}}
      {\olfileid{pl}{ntd}{rul}}

\olsection{قواعد و \usetoken{P}{derivation}}

\begin{explain}
قرار است دستگاه‌های استنتاج طبیعی با استدلال غیرصوریِ به‌کاررفته در
برهان‌های ریاضی تناظر نزدیکی داشته باشند (و ازاین‌رو تا حدی «طبیعی»
هستند). برهان‌های استنتاج طبیعی با فرض‌ها آغاز می‌شوند. سپس قواعد
استنتاج اعمال می‌شوند. فرض‌ها به‌وسیلهٔ قواعد استنتاج
\Intro{\lnot}، \Intro{\lif}، \iftag{FOL}{\Elim{\lor} و
  \Elim{\lexists}}{ و \Elim{\lor}} «!!{discharged}» می‌شوند و برای
روشنی، برچسب فرضِ !!{discharged} در کنار استنتاج قرار می‌گیرد.
\end{explain}

\begin{defn}[فرض]
\emph{فرض} هر !!{sentence}ای است که در بالاترین جایگاه یکی از شاخه‌ها
قرار گیرد.
\end{defn}

!!^{derivation}s در استنتاج طبیعی درخت‌های معینی از !!{sentence}s
هستند که در آن‌ها بالاترین !!{sentence}s فرض‌اند؛ و اگر
!!a{sentence} زیر یک، دو یا سه جملهٔ دیگر قرار گیرد، باید به‌درستی از
یک قاعدهٔ استنتاج نتیجه شود. !!{sentence}s بالای استنتاج را
\emph{مقدمات} و !!{sentence} زیر آن‌ها را \emph{نتیجهٔ} استنتاج
می‌نامند. قواعد به‌صورت جفت می‌آیند: برای هر !!{operator} یک قاعدهٔ
معرفی و یک قاعدهٔ حذف وجود دارد. این قواعد در نتیجه !!a{operator} را
معرفی می‌کنند یا !!a{operator} را از مقدمه‌ای از قاعده حذف می‌کنند.
برخی قواعد اجازه می‌دهند فرضی از نوعی معین \emph{!!{discharged}} شود.
برای نشان‌دادن اینکه کدام فرض به‌وسیلهٔ کدام استنتاج !!{discharged}
می‌شود، به فرض و استنتاج هر دو برچسب می‌دهیم. این امر با نوشتن فرض به
صورت «$\Discharge{!A}{n}$» نشان داده می‌شود.

% Only include this sentence is on of \land, lor, lif or lnot is defined.

\iftag{notprvNot,notprvAnd,notprvOr,notprvIf}{}%
	{مرسوم است که حتی اگر برخی از !!{operator}s تعریف شده باشند، قواعدی برای همهٔ آن‌ها، یعنی $\land$، $\lor$، $\lif$، $\lnot$ و $\lfalse$، در نظر بگیریم.}



\end{document}
